% ---------------------------------------------------------------------------
% Section: Phase 8 -- quantum consciousness, neuroscience and sequential evidence
% ---------------------------------------------------------------------------
\section{Phase 8: Quantum Consciousness, Neuroscience and Sequential Evidence}
\label{sec:phase8}

Phase~8 (20 October 2026) asked whether quantum theories of consciousness and the
neuroscience of decision making offer improvements to the inference speed or the
reasoning accuracy of \Sys. Each candidate idea had to pass four tests before any
code was written: it must be established rather than hypothesised; a classical,
CPU-only pipeline must inherit something from it beyond its classical description;
it must yield a result with an engineering consequence in the sense of
Section~\ref{sec:math}; and the consequence must be measurable on the reference
machine with the exact verifiers. The full review, with sources, is
\texttt{docs/research/quantum\_neuro\_review.md}; this section records the results.

\subsection{Quantum consciousness: a negative result}

Orchestrated objective reduction \cite{hameroff2014orchor} places consciousness
in quantum computations within microtubules. Tegmark's estimate of the
decoherence times of the relevant superpositions is $10^{-20}$--$10^{-13}$\,s
against dynamical timescales of $10^{-3}$--$10^{-1}$\,s \cite{tegmark2000decoherence};
the rebuttal of \cite{hagan2002quantum} lengthens the former by several orders of
magnitude without reaching the latter. Recent measurements of collective
ultraviolet emission from tryptophan networks in microtubules
\cite{babcock2024superradiance} show a modest quantum-yield enhancement and do
not demonstrate computation. Two results close the question for \Sys.

\begin{proposition}[Decoherence separation] \textnormal{[literature]}
\label{prop:q1}
With $\tau_{\mathrm{dec}}\le10^{-4}$\,s on the most favourable published estimate and
$\tau_{\mathrm{dyn}}\ge10^{-3}$\,s, any superposition of neural states has decohered into a
mixture over the interval of a decision. The description of cognition at that
timescale is a classical stochastic process.
\end{proposition}
\emph{Consequence.} The object is the one Section~\ref{sec:math} already models:
stochastic policies $\pi_1,\pi_2$ and a finite automaton over trace states.

\begin{proposition}[No transferable algorithm] \textnormal{[exact]}
\label{prop:q2}
A classical machine that follows the state of $n$ two-level systems stores $2^n$
complex amplitudes; in double precision the reference machine's 32\,GB hold
$n\le30$. No theorem separates BQP from BPP, so no speed-up is inherited by a
classical algorithm from the quantum character of a process.
\end{proposition}
\begin{proof}
$\dim\mathbb C^{2^n}=2^n$ and $2^{30}\times16$ bytes $=17$\,GB; the second claim is
the open status of the separation. \qed
\end{proof}
\emph{Consequence.} Even an established Orch~OR would contribute nothing computable to
a CPU inference loop (A-055). Quantum \emph{cognition}, the use of Hilbert-space
probability as a descriptive model of judgement
\cite{busemeyer2012quantum,pothos2013quantum}, has one structural lesson:
non-commuting measurements produce order effects, confirmed by the parameter-free
QQ equality on about seventy surveys \cite{wang2013qq,wang2014context} and
reported in language models \cite{qqllm2026}. Read for \Sys, an order effect is
anchoring of the deliberate answer on the draft, which the draft-blind slow path
(FR-023) removes and the coupling test (ER-004) measures
(Proposition~\ref{prop:collision}). The design needed no change.

\subsection{Neuroscience: the sequential form of the gate}

Three results of decision neuroscience have exact mathematical forms: evidence is
accumulated to a bound \cite{gold2007neural}, whose optimal form is Wald's
sequential probability ratio test \cite{wald1945sequential,wald1948optimum}; under
a cost of time or a deadline the optimal bound collapses
\cite{frazier2007deadline,drugowitsch2012cost}; and the quality of a confidence
signal is its accuracy--confidence association measured by ROC statistics
\cite{fleming2014metacognition}, with the all-or-none ignition of a global
workspace as the same threshold picture at system level
\cite{dehaene2011workspace}. Theorem~\ref{thm:gate} is the one-step instance;
three mechanisms of \Sys\ read a \emph{stream} of evidence and stop on a bound,
and each had a hand-set rule where a theorem exists. The language-model
literature reached the same rules independently
\cite{aggarwal2023adaptive,li2024esc,fu2025deepconf}.

\begin{proposition}[ROC dominance] \textnormal{[exact]}
\label{prop:n1}
If the ROC curve of a confidence signal $s'$ for predicting a wrong draft
dominates that of $s$ at every escalation rate, the best threshold gate on $s'$
has an objective \eqref{eq:objective} no larger than the best threshold gate on
$s$, for every $L$, $\lambda$ and $\Delta c$.
\end{proposition}
\begin{proof}
At a fixed escalation rate the compute term is equal and the error term is
$L$ times the wrong drafts not escalated, which dominance makes no larger;
minimising over the rate preserves the inequality. \qed
\end{proof}
\emph{Consequence.} The gate's signal is chosen by a stored ROC study against
exact labels (ER-021), not by convention.

\begin{proposition}[Exact early escalation on a running minimum] \textnormal{[exact]}
\label{prop:n2}
Let $w(n)$ be the lowest mean log-probability over complete windows of $W$
consecutive tokens among the first $n$ tokens of a draft. Then $w$ is
non-increasing in $n$, so for the rule ``escalate iff $w(N)<\theta$'' the
stopping time $T=\min\{n:w(n)<\theta\}$ gives exactly the full draft's decision,
and no rule on the same stream escalates such a draft earlier without a
non-zero disagreement rate.
\end{proposition}
\begin{proof}
A minimum over a growing set; $w(N)\le w(T)<\theta$; before $T$ every window is
$\ge\theta$. \qed
\end{proof}
\emph{Consequence.} Early escalation on the lowest-window signal needs no margin,
minimum token count or agreement study (FR-083). It is a decision, not a
prediction.

\begin{proposition}[Wald's bound on the mean rule] \textnormal{[exact under i.i.d.\ Gaussian tokens]}
\label{prop:n3}
Any sequential test of $\mu=\tau$ against $\mu=\tau-\delta$ on tokens with
variance $\sigma^2$ and errors $(\alpha,\beta)$ has
$\mathbb E_1[N]\ge\big[(1-\beta)\ln\frac{1-\beta}{\alpha}+\beta\ln\frac{\beta}{1-\alpha}\big]\big/\frac{\delta^2}{2\sigma^2}$.
\end{proposition}
\emph{Measured.} On the 40 open-ended drafts of the early-escalation study the
per-token spread is $\sigma\approx0.28$ and the margin is $\delta=0.04$: the bound
is 260 tokens at $\alpha=\beta=0.05$, against drafts of 47--291 tokens (median
133), and the first tokens of a draft are systematically more confident than the
rest. \emph{Consequence.} The mean-based early escalation of P7-T03 is an
empirical heuristic that cannot be made more aggressive by a sequential test; it
is kept at its measured setting and labelled as such (A-056).

\begin{proposition}[Posterior stopping with a collapsing bound] \textnormal{[exact for the rule; strong for optimality]}
\label{prop:n4}
With $k$ probed answers clustered into categories, counts $v$ and a uniform
Dirichlet prior, the posterior probability that the plurality answer is the
majority answer is
$1-I_{1/2}(v_{\mathrm{top}}+1,\,k-v_{\mathrm{top}}+1)
=1-2^{-k-1}\sum_{j=v_{\mathrm{top}}+1}^{k+1}\binom{k+1}{j}$.
Under a per-probe cost and a hard budget the Bayes-optimal rule stops at the
first $k$ with this probability at or above a bound $\theta(t_k)$ that is
non-increasing in the budget used \cite{frazier2007deadline}.
\end{proposition}
\emph{Values.} One probe 0.75; two agreeing 0.875; $[2,1]$ 0.6875; $[3,1]$ 0.8125.
With $\theta$ falling linearly from 0.85 to 0.70 the rule reproduces the
``two consecutive agreeing probes'' rule early in the budget, never stops on one
probe, votes over all probes (Theorem~\ref{thm:plurality}) rather than the last
two, and lets a $[3,1]$ split stop once about 45\,\% of the budget is used
(FR-084). Verifiable tasks are untouched: the exact verifier is the stopping
signal (Theorem~\ref{thm:precision}).

\begin{proposition}[Optimal probe interval] \textnormal{[exact under an exponential model]}
\label{prop:n5}
If the reasoning length at which the answer becomes available is exponential
with mean $m$ tokens, a probe costs $c_p$ seconds and decoding runs at $v$
tokens per second, the periodic interval minimising probe cost plus overshoot is
$\Delta^\star=\sqrt{2c_pvm}$, with $\mathrm{cost}(\Delta)/\mathrm{cost}(\Delta^\star)
=\tfrac12(\Delta/\Delta^\star+\Delta^\star/\Delta)$.
\end{proposition}
\emph{Measured.} Idle: $c_p\approx0.7$\,s, $v\approx8$, $m\approx300$ give
$\Delta^\star\approx58$ tokens; the configured gap of 64 costs 0.5\,\% more than
the optimum (1\,\% under the shared load). The schedule is validated and
unchanged (A-059). This is the quantitative content of the ``quantum Zeno''
intuition that measuring often is costly; here a probe has no back-action, only
a price.

\subsection{Measurements}
\label{sec:phase8meas}

All measurements were taken on the shared reference machine (load average 5--11,
other sessions compiling); accuracy and log-probabilities do not depend on load,
wall-clock figures do.

\textbf{Confidence-signal study (ER-021).} 150 GSM8K test items were drafted by
System~1 under the production prompt and sampling (temperature 0.2, top-$p$ 0.8,
no thinking) and scored by the exact numeric verifier: 31 drafts are wrong
(accuracy 0.793). Table~\ref{tab:signals} gives, for each candidate signal, the
AUROC for predicting a wrong draft (lower signal flags the draft; 400-resample
bootstrap interval) and the recall and precision when the least confident 10\,\%
and 20\,\% of drafts are escalated.

\begin{table}[t]
\caption{Which confidence signal of a System~1 draft predicts a wrong draft (150 GSM8K drafts, 31 wrong)}
\label{tab:signals}
\centering\small
\begin{tabular}{@{}lccc@{}}
\toprule
Signal & AUROC [95\,\% CI] & recall / precision at 10\,\% & at 20\,\% \\
\midrule
lowest token log-prob & \textbf{0.646} [0.539, 0.746] & 0.23 / 0.47 & 0.35 / 0.37 \\
mean top-1/top-2 margin & 0.614 [0.493, 0.729] & 0.23 / 0.47 & 0.39 / 0.40 \\
lowest 16-token window of the margin & 0.611 [0.510, 0.715] & 0.06 / 0.13 & 0.26 / 0.27 \\
lowest 16-token window log-prob \cite{fu2025deepconf} & 0.581 [0.488, 0.698] & 0.13 / 0.27 & 0.29 / 0.30 \\
negative mean top-5 entropy & 0.573 [0.453, 0.679] & 0.16 / 0.33 & 0.26 / 0.27 \\
mean log-prob (the fixed rule) & 0.564 [0.449, 0.672] & 0.13 / 0.27 & 0.26 / 0.27 \\
bottom-10\,\% tokens & 0.555 [0.459, 0.675] & 0.10 / 0.20 & 0.23 / 0.23 \\
tail (last 32 tokens) & 0.492 [0.366, 0.631] & 0.19 / 0.40 & 0.23 / 0.23 \\
\bottomrule
\end{tabular}
\end{table}

Three facts follow. First, \emph{the fixed rule is blind on mathematics}: at the
production threshold ($-0.175$ at the default $\lambda$) the mean-log-prob rule
escalates 0 of 150 drafts and catches 0 of 31 wrong ones, because mathematics
drafts have a mean log-prob of about $-0.03$ whether right or wrong; on the
unverified mathematics path every wrong draft was kept, and because such drafts
never enter the uncertainty band the agreement samples were never drawn either.
Second, \emph{token confidence is a weak metacognitive signal} for this model on
this task: the best signal is the single lowest token, and the window signal's
interval includes chance, so no ROC dominates (Proposition~\ref{prop:n1}) and
the signal that Proposition~\ref{prop:n2} makes exact to act on is not the one
worth acting on. Third, \emph{Theorem~\ref{thm:gate} forbids escalating on any of
them at the default price}: escalation pays only if
$\Pr[\text{wrong}\mid\text{escalated}]>\hat e_2+\lambda\Delta c/(L\cdot\text{cost unit})
=0.15+0.2\cdot150/60=0.65$, and the highest precision reached at an escalation
rate of at least 10\,\% is 0.47. By Proposition~\ref{prop:cascade} the expected
accuracy of a gate escalating the 20\,\% least confident drafts by the lowest
token is 0.847--0.867 (System~2 accuracy 0.9--1.0) against 0.793, for about 30\,s
of mean latency per question, which the default price rates as a loss.
\emph{Decision:} the window signal is recorded on every decision and read by the
calibrated gate, but its fixed-rule threshold stays off (A-057): the mechanism is
exact, the measured signal is not strong enough to use it. The repository keeps
this negative result next to the positive ones.

\textbf{Agreement with cheap samples.} The gate's second signal, agreement between
the draft and $k=2$ cheap System~1 samples, was never computed on mathematics
because the samples are drawn only inside the uncertainty band of the mean
log-prob. For the first 60 drafts the samples were drawn exactly as production
draws them and compared with the draft's final number (12 of 60 drafts wrong;
54\,s per item on the loaded machine). Agreement has AUROC 0.759 [0.618, 0.901]
and, with ties broken by the lowest token, 0.823 [0.690, 0.940], against 0.720
for the lowest token and 0.611 for the window on the same drafts. The rule
``any disagreement'' (the configured threshold at the default care) escalates
25\,\% of drafts, catches 8 of the 12 wrong ones (precision 0.53) and raises the
expected cascade accuracy from 0.800 to 0.908--0.933. This is what
Theorem~\ref{thm:plurality} predicts: a wrong low-temperature draft is usually
one whose answer holds no majority at temperature 0.8. \emph{Decision
(implemented, FR-085):} unverified mathematics drafts always get their two
samples (\texttt{gate.samples\_on\_unverified\_domains}), whatever their token
confidence says (A-060). Its precision is still below the 0.65 bar of
Theorem~\ref{thm:gate} at the default price, so the calibrated gate, once
promoted, will price it rather than apply it unconditionally.

\textbf{Exit rules on factual prompts (FR-084).} A paired Lab run on six factual
prompts (run \texttt{73cd5bba}, 512-token budget, shared machine) compared the
consistency rule with the posterior rule. Both gave the same factual content on
every prompt; reasoning tokens were 357 against 382 on average (wall 93\,s
against 101\,s), two prompts exhausting the budget under both rules. Where they
diverged most, the probes disagreed first and then agreed: the consistency rule
stopped at the two consecutive agreeing probes (183 tokens) while the posterior
rule saw a $[2,1]$ split below its bound and continued to $[3,1]$ (318 tokens),
as Proposition~\ref{prop:n4} prescribes. That caution costs tokens on this set,
so by the criterion fixed in advance (A-058) the posterior rule stays an option
and the consistency rule remains the default; the comparison is indicative
($n=6$, no exact labels).


\subsection{What changed}

The crate \texttt{nafs-core} gains a \texttt{sequential} module (the window signal and
its streaming tracker, the probe posterior with the exact beta tail, the
collapsing bound, Wald's bound and the probe interval, each with tests of the
property its proposition states). Every gate decision now records
\texttt{low\_window\_logprob}; the calibrated gate reads it as an eighth feature; the
fixed rule applies it for the measured domains and stops the draft at the
crossing; the adaptive exit offers the posterior rule per request and in the Lab
(\texttt{slow\_exit\_posterior} against \texttt{slow\_exit\_consistency}). Nothing was
added for the quantum-consciousness ideas; Propositions~\ref{prop:q1}--\ref{prop:q2}
are the record of why.
